\documentclass[10pt,a4paper]{article}
\usepackage[T1]{fontenc}
\usepackage[utf8]{inputenc}
\usepackage{lmodern}
\usepackage[french]{babel}
\usepackage[margin=1.9cm]{geometry}
\usepackage{amsmath}
\usepackage{graphicx}
\usepackage{subcaption}
\usepackage{booktabs}
\usepackage{siunitx}
\usepackage{hyperref}
\usepackage{float}
\sisetup{locale=FR}
\graphicspath{{../resultats/}{resultats/}{../photos/}{photos/}}
% \IfFileExists ne consulte pas \graphicspath : mêmes dossiers pour \photo.
\makeatletter
\def\input@path{{../resultats/}{resultats/}{../photos/}{photos/}}
\makeatother
\hypersetup{hidelinks}
\setlength{\parindent}{0pt}
\setlength{\parskip}{3pt}
\setlength{\textfloatsep}{8pt}
\setlength{\intextsep}{6pt}

% Emplacement visible si l'image n'a pas encore été ajoutée.
\newcommand{\photo}[2]{%
  \IfFileExists{#1}{\includegraphics[width=#2]{#1}}{%
    \fbox{\parbox[c][2.7cm][c]{#2}{\centering\footnotesize Image à ajouter :\\\texttt{\detokenize{#1}}}}%
  }%
}

\title{\vspace{-1cm}\textbf{Calibration d'une caméra}\\\large Séquence Vidéo -- TP 1}
\author{Roméo Agostino -- Master 2 VMI\\\small Code : \url{https://github.com/RomeoAg13/Sequence-video}}
\date{}
\begin{document}
\maketitle
\vspace{-0.9cm}

\section{Introduction}
L'objectif de ce travail pratique est de réaliser la calibration d'une caméra à partir de plusieurs photographies d'une mire, pour estimer les paramètres intrinsèques de la caméra, caractériser la distorsion optique et déterminer la position et l'orientation d'un objet par rapport à la caméra.

Pour ce TP, j'ai utilisé la caméra de mon téléphone. La mire est un damier composé de $9\times7$ cases, soit $8\times6=48$ coins intérieurs détectables.

N'ayant pas de damier, j'ai affiché la mire sur l'écran d'un ordinateur portable. J'ai réalisé 8 photographies avec différentes positions et orientations de la caméra, pour varier l'observation.

Les images utilisées ont une résolution de $1536\times1152$ pixels et la taille d'une case de la mire a été estimée à environ \SI{2}{\centi\metre}.

\section{Méthode}
Les images sont chargées avec \texttt{cv2.imread()}, converties en niveaux de gris, puis les 48 coins intérieurs sont détectés avec \texttt{findChessboardCorners()} et affinés avec \texttt{cornerSubPix()}.

Pour chaque photographie, \texttt{points\_3d} contient les coordonnées des coins sur le damier et \texttt{points\_2d} leurs coordonnées en pixels. Les 8 photographies ont été utilisées pour la calibration avec \texttt{calibrateCamera()}.

\begin{figure}[H]
  \centering
  \photo{coins_gris_1.jpg}{0.47\linewidth}\hfill
  \photo{coins_couleur_1.jpg}{0.47\linewidth}\\[4pt]
  \photo{coins_gris_2.jpg}{0.47\linewidth}\hfill
  \photo{coins_couleur_2.jpg}{0.47\linewidth}
  \caption{Détection des 48 coins intérieurs sur plusieurs photographies}
\end{figure}

\section{Résultats de calibration}
La matrice intrinsèque estimée est :
\[
K=\begin{pmatrix}
1104{,}62 & 0 & 769{,}64\\
0 & 1104{,}24 & 580{,}68\\
0 & 0 & 1
\end{pmatrix}.
\]
Les focales $f_x$ et $f_y$ sont exprimées en pixels. Les deux focales sont très proches et le point principal est voisin du centre de l'image.

\begin{table}[H]
\centering\small
\begin{tabular}{ll}
\toprule
Paramètre & Valeur \\
\midrule
Erreur RMS de reprojection & $0{,}177$ pixel \\
Distorsion radiale $k_1,k_2,k_3$ & $0{,}1783$ , $-0{,}3943$ , $-0{,}6383$ \\
Distorsion tangentielle $p_1,p_2$ & $0{,}000298$ , $-0{,}001165$ \\
Dimensions de capteur utilisées (estimées) & $9{,}80 \times 7{,}30$ mm \\
Focale physique estimée & $7{,}048$ mm $\approx \boxed{0{,}00705\ \mathrm{m}}$ \\
Champ de vision horizontal / vertical & $69{,}62^\circ$ / $55{,}09^\circ$ \\
\bottomrule
\end{tabular}
\caption{Valeurs obtenues avec \texttt{calibrateCamera()} et \texttt{calibrationMatrixValues()}.}
\end{table}

L'erreur RMS de $0{,}177$ pixel indique un faible écart moyen entre les coins détectés et les points reprojetés. La focale en mètres a été estimée avec \texttt{calibrationMatrixValues()} : elle dépend des dimensions de capteur supposées, et n'est donc pas une mesure physique exacte.

\section{Correction de la distorsion}
La fonction \texttt{undistort()} permet de corriger les déformations optiques à partir des paramètres calculés. Elle ne supprime pas l'effet de perspective dû à l'inclinaison du damier.

\begin{figure}[H]
\centering
\begin{subfigure}{0.39\linewidth}
  \centering\photo{photo_8.jpg}{\linewidth}
  \caption{Image originale}
\end{subfigure}\hspace{0.03\linewidth}
\begin{subfigure}{0.39\linewidth}
  \centering\photo{corrigee_photo_8.jpg}{\linewidth}
  \caption{Après correction.}
\end{subfigure}\\[4pt]
\begin{subfigure}{0.39\linewidth}
  \centering\photo{photo_5.jpg}{\linewidth}
  \caption{Image originale}
\end{subfigure}\hspace{0.03\linewidth}
\begin{subfigure}{0.39\linewidth}
  \centering\photo{corrigee_photo_5.jpg}{\linewidth}
  \caption{Après correction.}
\end{subfigure}
\caption{Comparaison avant/après la correction de distorsion.}
\end{figure}

\section{Estimation de pose et cube 3D}
\texttt{solvePnP()} utilise les correspondances entre coins 3D et 2D pour estimer la rotation (\texttt{rvec}) et la translation (\texttt{tvec}) du damier dans le repère de la caméra. Les axes sont affichés avec \texttt{drawFrameAxes()}. L'origine (0,0,0) est placée sur le premier coin intérieur du damier.

Rotation (rvec) :
\[
rvec=\begin{pmatrix}
0{,}289\\
0{,}864\\
0{,}470
\end{pmatrix}
\]

Translation (tvec) :
\[
tvec=\begin{pmatrix}
-1{,}851\\
-3{,}987\\
15{,}256
\end{pmatrix}
\]

Pour aller un peu plus loin, j'ai construit les 8 sommets d'un cube en 3D. La fonction \texttt{projectPoints()} projette ces sommets sur la photographie, puis \texttt{line()} permet de dessiner les douze arêtes en couleur.

\begin{figure}[H]
  \centering
  \photo{axes_cube.jpg}{0.62\linewidth}
  \caption{Axes du repère et cube 3D projeté sur la mire.}
\end{figure}
\end{document}